\section{Scoring}
\label{sec:scoring}

After a rollout, a scenario is scored on four CLEAR-aligned dimensions: efficacy,
cost, reliability, and latency. Efficacy is the only dimension involving judges,
and even there a third is deterministic.

\subsection{Efficacy: a hybrid, partly judge-independent score}

For a v0.2 scenario carrying a ground-truth world, efficacy is
\begin{equation}
\text{efficacy} = 0.4 \cdot \text{task\_completion} + 0.3 \cdot \text{tool\_selection} + 0.3 \cdot \text{state\_verification},
\end{equation}
where the first two terms are LLM-judge consensus and the third is deterministic.
A full third of efficacy is therefore objective. For legacy scenarios with no
ground truth, state verification is inapplicable and efficacy renormalizes to an
equal split over the two judge dimensions; all v1 public scenarios carry ground
truth, so the hybrid form is the operative one.

\paragraph{State verification (deterministic).} The final world, after the
conversation, is checked against the scenario's \texttt{expected\_state\_changes}
with no model in the loop. Each assertion uses the tiny vocabulary of
Section~\ref{sec:design} (\texttt{equals}; \texttt{increased\_by}/%
\texttt{decreased\_by} with float tolerance 0.01; \texttt{contains}), and the state
score is the fraction of assertions that pass. An empty assertion list scores 1.0
if and only if the world is unchanged, the no-unauthorized-mutation contract. This
component is free to compute and not gameable by fluent text.

\paragraph{Task completion (judge).} Each judge grades whether the agent
accomplished the user's goals on a COMPLETE (1.0) / PARTIAL (0.5) / FAILED (0.0)
anchor, with $\pm0.1$ adjustments for clarifying questions, error recovery, scope
awareness, and dependency handling. The full rubric is published as code in
\texttt{eval/scoring/rubrics.py}.

\paragraph{Tool selection (judge).} Each judge grades every tool call on five
dimensions: selection correctness, parameter accuracy, sequencing, necessity, and
omissions.

\paragraph{Combined judge prompt.} Task completion and tool selection are scored
in a single call per judge. The combined prompt presents the shared context and the
transcript once and asks for both dimensions in one JSON response nesting the two
per-rubric shapes. The per-dimension criteria are the same published rubric text,
reorganized under one shared context rather than rewritten, so the consensus math,
per-judge columns, artifacts, aggregation, and downstream checks are unchanged. The
benefit is that the judge call count is halved and, because the transcript is no
longer duplicated across two calls, judge input tokens drop by roughly 45\%. The
honest tradeoff is failure coupling: because both dimensions ride one API call and
one parsed body, an API failure drops the judge from both panels, and a parse
failure (unrecoverable JSON after one retry, or a body in which either dimension is
missing or has an invalid score) flags the judge as parse-failed for both. The
legacy two-call path remains available behind a flag for A/B validation.

\subsection{The judge panel and consensus}

Three judges score every transcript independently: two open-weight models from
different labs (Kimi K2.6 from Moonshot and GLM-4.6 from Zhipu, both served via
OpenRouter) and one frontier reference (Claude Opus, via the Anthropic API). Two
principles govern the panel. First, \textbf{no judge is also a model under test};
an earlier panel violated this and was corrected. Second, \textbf{lab diversity},
so a disagreement signals genuine ambiguity rather than a shared blind spot. Open
judges run over hosted APIs, so the full panel needs no GPU.

\paragraph{Median consensus.} Consensus is the median of the valid judge scores.
For an $n=3$ panel the median is the middle judge, so a single rogue or
leniency-drifted judge cannot drag the consensus the way a mean would; for $n=2$ the
median equals the mean. Bradley-Terry was rejected because it models pairwise
preference data, while COT Bench produces absolute pointwise rubric scores. Every
individual judge score is published, so any other aggregation can be recomputed.

\paragraph{Judge-failure handling.} A parse failure is not treated as a real 0.0,
which would crater both the consensus and the agreement rate indistinguishably from
a genuine low score. The judge is called once more; if parsing still fails, the
judge is excluded from consensus, agreement, and max-disagreement but retained in
the per-row record. An API failure drops the judge from that row. Because failures
shrink the panel, every row publishes explicit accounting: the number of valid
judges per dimension, parse-failure and API-failure counts, and a \texttt{degraded}
flag when fewer than two valid judges remain. Agreement is reported as null with
fewer than two valid judges, since a single grader is not perfect agreement.

\subsection{Per-judge accountability}

For every contestant and every judge, the board publishes that judge's mean task
completion and tool selection for the model and its delta against the full-panel
consensus (\texttt{judge\_deltas}; $\text{delta} = \text{consensus\_mean} -
\text{judge\_mean}$, so a positive delta means the judge rates the model lower than
its peers). This makes any judge's systematic favoritism toward a model a published
number. The one remaining same-lab pairing (the Claude judge shares a lab with the
Claude contestants) is handled as a special case: each Anthropic contestant carries
a \texttt{same\_lab\_check} block reporting its means over the open judges only and
the delta versus the full panel, so any same-lab inflation is measured, not
assumed away. Anyone can recompute consensus with the same-lab judge excluded from
the published per-judge scores.

\subsection{Cost, reliability, and latency}

\paragraph{Cost.} $\text{cost\_per\_task} = \text{in\_tokens} \cdot p_{\text{in}} +
\text{out\_tokens} \cdot p_{\text{out}}$ at published per-token provider pricing,
tracked in \texttt{eval/config.py} and dated per run.

\paragraph{Latency.} Wall-clock time summed across the agent's turns only; the user
and tool simulators are excluded, so the number reflects the wait a user would
actually experience.

\paragraph{Reliability and $\text{pass}^k$.} Each scenario runs three times. The
board reports the pass rate (fraction above a 0.7 threshold), consistency (1.0
minus the high-low spread), variance, and $\text{pass}^k$. Borrowed from
tau-bench~\cite{taubench2024}, $\text{pass}^k$ is the probability that \emph{every
one} of $k$ independent trials succeeds, the opposite of $\text{pass}@k$. For
i.i.d. trials it decays as $p^k$, so a model that passes 2 of 3 runs has
$\text{pass}^1 = 0.67$ but $\text{pass}^3 = 0$: one failure across repeats sinks
it. It is estimated per scenario as the average over all size-$k$ subsets of the
$n$ collected trials of the all-pass indicator, with closed form
$\binom{c}{k}/\binom{n}{k}$ for $c$ passing runs out of $n$. Each model publishes
$\text{pass}^k$ alongside, not replacing, the pass@3 reliability. With three runs,
$k \in \{1,2,3\}$; more repeats extend the autonomy horizon.

\paragraph{Failure-mode taxonomy.} Every run below the 0.7 pass threshold is
additionally classified into one of six failure modes (tool-selection-error,
policy-violation, incomplete-task, premature-end, wrong-parameters,
hallucinated-capability), deterministic-first: state-grader evidence and the
premature-end instrumentation take precedence over keyword matching on the
judges' already-written reasoning, with no additional LLM calls; each model
publishes a per-mode failure profile, the diagnostic practitioners act on.

\subsection{CLEAR composite}

The composite adopts four of CLEAR's five dimensions~\cite{clear2025}: Cost,
Latency, Efficacy, and Reliability. CLEAR's fifth dimension, Assurance
(safety and policy adherence), is not scored as a separate axis here; policy
behavior is instead exercised inside the scenarios themselves (the scope-management
and adversarial-input categories) and graded through the same efficacy machinery.
Dimensions are min-max normalized across the evaluated models, then weighted:
\begin{equation}
\text{CLEAR} = 0.35\,\text{eff}_{\text{norm}} + 0.25\,\text{rel}_{\text{norm}} + 0.20\,(1 - \text{cost}_{\text{norm}}) + 0.20\,(1 - \text{lat}_{\text{norm}}).
\end{equation}
Cost and latency are inverted because lower is better. Because normalization is
field-relative, adding or removing a model rescales every other model's composite,
so CLEAR scores are only comparable within a run; cross-run comparison must use raw
per-dimension scores. With a single model, CLEAR equals raw efficacy.

\subsection{Anti-gaming: the do-nothing agent}

The RDI audit showed trivial agents can game several agentic benchmarks to
near-perfect scores~\cite{rdi2026}. COT Bench includes a deterministic do-nothing
agent that makes no tool calls and returns a single trivial deflecting message. It
is expected to score near zero on both halves of efficacy: the state checks go to
0.0 because the world is never mutated (judge-independent and not gameable by
text), and the judges floor because the transcript shows no tools used and no
goals resolved. The lone case where doing nothing is correct is a
no-unauthorized-mutation scenario with an empty assertion list, where the
do-nothing agent ``passes'' precisely because it changed nothing, which is the
intended contract. The probe is deliberately excluded from the published board and
from the bootstrap and rank-band math; publishing the result of a real judged run
is a methodology credential that a near-zero floor exists.

\subsection{Human calibration (protocol defined; study pending)}

Inter-judge agreement is necessary but not sufficient: three judges can agree and
all be wrong, the MT-Bench and PoLL precedent~\cite{mtbench2023,poll2024}. The
calibration protocol is defined and tooled now; the study runs after the rehearsal
run produces transcripts (it needs human hours, not API spend). From a run's
artifacts, 50--80 transcripts (default 60) are drawn stratified by domain
$\times$ category $\times$ difficulty band, one per (scenario, model) to avoid
pseudo-replication, with the private holdout excluded by default. The sample is
emitted as a blind labeling workbook (opaque token per sheet, rubric anchors
printed inline, scores and model identity held in a separate key file). Filled
labels are scored for human-versus-judge Krippendorff's $\alpha$, mean absolute
difference, and Pearson correlation, against both the consensus and each individual
judge. A materially higher human-vs-consensus $\alpha$ than the worst
human-vs-single-judge $\alpha$ would be direct evidence the panel beats any one
judge, the PoLL result~\cite{poll2024}. The published number will be labeled a
calibration smoke test, not a powered study.
